\begin{ouverture}[geometrie]{15}{Le théorème de Pythagore}{Partie 4 · Géométrie et espace}
{Triangle rectangle · Longueurs · Réciproque · Contraposée}
\objectif{Calculer une longueur dans un triangle rectangle.}
\objectif{Démontrer qu’un triangle est rectangle ou qu’il ne l’est pas.}
\objectif{Distinguer théorème, réciproque et contraposée.}
\prerequis{Carrés, racines carrées et raisonnement (chapitres 4 et 12).}
\end{ouverture}
\pagestyle{otMain}\essentiel\fontsize{11}{15}\selectfont
\phantomsection\label{cours-15-1}
\essSection{1}{Calculer dans un triangle rectangle}
\begin{theoreme}[Théorème de Pythagore]Si un triangle ABC est rectangle en A, alors $BC^2=AB^2+AC^2$. Le côté opposé à l’angle droit, $[BC]$, est l’hypoténuse : c’est le plus long côté. Le théorème relie les aires des carrés construits sur les trois côtés.\end{theoreme}
\label{fig-cours-15-1-cours-pythagore}\figureGeo{cours-pythagore}

\begin{exemple}Si $AB=3$ et $AC=4$ cm, alors $BC^2=9+16=25$, donc $BC=\sqrt{25}=5$ cm. Si $BC=13$ et $AB=5$ cm, alors $AC^2=13^2-5^2=144$, donc $AC=12$ cm.\end{exemple}
\begin{attention}On additionne les carrés pour chercher l’hypoténuse ; on soustrait pour chercher un côté de l’angle droit. Écrire d’abord l’égalité avec les noms des points aide à éviter cette confusion. Une longueur est positive : on prend la racine carrée.\end{attention}
\clearpage
\phantomsection\label{cours-15-2}
\essSection{2}{Réciproque et contraposée}
\begin{theoreme}[Réciproque du théorème de Pythagore]Si, dans un triangle ABC, $BC^2=AB^2+AC^2$, alors ce triangle est rectangle en A.\end{theoreme}
\begin{propriete}[Contraposée]Si, dans un triangle ABC, $BC^2\ne AB^2+AC^2$, alors ce triangle n’est pas rectangle en A. Lorsque $[BC]$ est le plus long côté, un angle droit ne pourrait être qu’en A, face à ce côté : le triangle n’est alors pas rectangle.\end{propriete}
\begin{savoirfaire}[Prouver ou réfuter]Identifier le plus long côté. Calculer séparément son carré et la somme des carrés des deux autres côtés. En cas d’égalité, la réciproque prouve que le triangle est rectangle ; sinon, la contraposée prouve qu’il ne l’est pas.\end{savoirfaire}
\begin{exemple}Côtés 6, 8 et 10 : $10^2=100=36+64$, le triangle est rectangle. Côtés 6, 8 et 11 : $11^2=121\ne100$, il ne l’est pas. Un dessin qui paraît rectangle ne suffit pas.\end{exemple}
\begin{savoirfaire}[Rédiger]Nommer le triangle et son angle droit, citer le théorème, écrire l’égalité, calculer, prendre la racine positive et conclure avec l’unité. Garder la valeur exacte avant l’arrondi final.\end{savoirfaire}

\clearpage\rubrique{Découverte}
\begin{blocLivre}{ex-15-1}\begin{exercice}[1]{Chercher l’hypoténuse}
ABC est rectangle en A, $AB=6$ cm et $AC=8$ cm. Calculer BC.
\label{fig-ex-15-1}\figureGeo{pyth-01}
\end{exercice}\end{blocLivre}
\begin{blocLivre}{ex-15-2}\begin{exercice}[1]{Chercher un côté de l’angle droit}
ABC est rectangle en A, $AB=5$ cm et $BC=13$ cm. Calculer AC.
\label{fig-ex-15-2}\figureGeo{pyth-02}
\end{exercice}\end{blocLivre}
\begin{blocLivre}{ex-15-3}\begin{exercice}[1]{Une racine non entière}
ABC est rectangle en A avec $AB=7$ cm et $AC=9$ cm. Donner BC exactement, puis au dixième.
\label{fig-ex-15-3}\figureGeo{pyth-03}
\end{exercice}\end{blocLivre}
\begin{blocLivre}{ex-15-4}\begin{exercice}[1]{L’hypoténuse donnée}
ABC est rectangle en A, $BC=17$ cm et $AB=15$ cm. Calculer AC.
\label{fig-ex-15-4}\figureGeo{pyth-04}
\end{exercice}\end{blocLivre}
\begin{blocLivre}{ex-15-5}\begin{exercice}[1]{Une réciproque}
Un triangle ABC a pour côtés $AB=9$, $AC=12$ et $BC=15$ cm. Est-il rectangle ?
\label{fig-ex-15-5}\figureGeo{pyth-05}
\end{exercice}\end{blocLivre}

\rubrique{Entraînement}
\begin{blocLivre}{ex-15-6}\begin{exercice}[2]{La contraposée}
Un triangle ABC a pour côtés $AB=4$, $AC=5$ et $BC=7$ cm. Est-il rectangle ?
\label{fig-ex-15-6}\figureGeo{pyth-06}
\end{exercice}\end{blocLivre}
\begin{blocLivre}{ex-15-7}\begin{exercice}[2]{Diagonale d’un rectangle}
Un rectangle mesure 8 cm sur 6 cm. Calculer la longueur de sa diagonale.
\label{fig-ex-15-7}\figureGeo{pyth-07}
\end{exercice}\end{blocLivre}
\begin{blocLivre}{ex-15-8}\begin{exercice}[2]{Une échelle contre un mur}
Une échelle de 2,5 m repose contre un mur perpendiculaire au sol. Son pied est à 0,8 m du mur. À quelle hauteur arrive-t-elle, au centimètre près ?
\label{fig-ex-15-8}\figureGeo{pyth-08}
\end{exercice}\end{blocLivre}
\begin{blocLivre}{ex-15-9}\begin{exercice}[2]{Une distance dans un repère}
Dans un repère orthonormé d’unité 1 cm, on donne $A(0\,;0)$, $B(6\,;0)$ et $C(0\,;8)$. Calculer BC.
\label{fig-ex-15-9}\figureGeo{pyth-09}
\end{exercice}\end{blocLivre}
\begin{blocLivre}{ex-15-10}\begin{exercice}[2]{Trouver le bon sommet}
$AB=10{,}5$ cm, $AC=6{,}3$ cm et $BC=8{,}4$ cm. Démontrer que ABC est rectangle et préciser le sommet de l’angle droit.
\label{fig-ex-15-10}\figureGeo{pyth-10}
\end{exercice}\end{blocLivre}

\rubrique{Problèmes}
\begin{blocLivre}{ex-15-11}\begin{exercice}[3]{Un losange}
Un losange a des diagonales de 10 cm et 24 cm, qui se coupent perpendiculairement en leur milieu. Calculer son côté et son périmètre.
\label{fig-ex-15-11}\figureGeo{pyth-11}
\end{exercice}\end{blocLivre}
\begin{blocLivre}{ex-15-12}\begin{exercice}[3]{Retrouver une aire}
Un rectangle a une diagonale de 13 cm et une largeur de 5 cm. Calculer sa longueur puis son aire.
\label{fig-ex-15-12}\figureGeo{pyth-12}
\end{exercice}\end{blocLivre}
\begin{blocLivre}{ex-15-13}\begin{exercice}[3]{Deux diagonales}
Un câble suit la diagonale d’un rectangle de 8 m sur 6 m, puis celle d’un rectangle de 10 m sur 24 m. Calculer la longueur totale du câble.
\label{fig-ex-15-13}\figureGeo{pyth-13}
\end{exercice}\end{blocLivre}
\begin{blocLivre}{ex-15-14}\begin{exercice}[3]{Une décision exacte}
Une barre rigide de 8 cm doit relier les extrémités de deux côtés perpendiculaires de 4 cm et 7 cm. Est-elle assez longue ? Justifier sans arrondi.
\label{fig-ex-15-14}\figureGeo{pyth-14}
\end{exercice}\end{blocLivre}
\begin{blocLivre}{ex-15-15}\begin{exercice}[3]{Choisir un trajet}
Un terrain rectangulaire mesure 12 m sur 5 m. Comparer un trajet le long de deux côtés consécutifs avec le trajet suivant la diagonale. Quelle longueur gagne-t-on avec la diagonale ?
\label{fig-ex-15-15}\figureGeo{pyth-15}
\end{exercice}\end{blocLivre}

\clearpage
